\section{De K1.5 a K2: la evolución técnica de Moonshot AI (月之暗面)}

\subsection{K1.5: validación de la técnica de aprendizaje por refuerzo}

K1.5 fue principalmente la validación de la ruta técnica de aprendizaje por refuerzo: Moonshot AI fue uno de los equipos que invirtieron en esta ruta desde etapas tempranas. Hallazgo clave:

\begin{importantbox}{El Reward de extremo a extremo supera al Process Reward}
Hallazgo técnico importante de K1.5: no es muy necesario un Process Reward Model (PRM) o una Value Function. Usar directamente un Reward de extremo a extremo permite lograr un entrenamiento muy bueno. El Process Reward incluso puede tener efectos secundarios durante el entrenamiento. Al principio, esta no era una conclusión muy clara.
\end{importantbox}

\subsection{K2: en busca de un mejor modelo base}

K2 tiene dos direcciones centrales en sus objetivos de diseño:

\subsubsection{Dirección uno: Token Efficiency, maximizar el valor de cada porción de datos}

\begin{importantbox}{Por qué Token Efficiency es más importante que Compute Efficiency}
El crecimiento de los datos de alta calidad es muy lento: puede aproximarse como una constante (por ejemplo, 30T Tokens de alta calidad). Bajo esta restricción:
\begin{itemize}
    \item \textbf{Compute Efficiency} (entrenar más rápido): aunque tiene valor, no eleva el límite superior de la inteligencia: solo se termina de entrenar más rápido con los mismos datos
    \item \textbf{Token Efficiency} (aprender mejor): con los mismos datos, el modelo absorbe más, la tasa de compresión es mayor y el loss baja más rápido: eleva directamente el límite superior de la inteligencia
\end{itemize}
\end{importantbox}

Moonshot AI mejora el Token Efficiency de tres maneras:

\begin{enumerate}
    \item \textbf{Optimizador Muon}: un nuevo optimizador basado en la ortogonalización de matrices, que sustituye a Adam, usado durante 10 años. En condiciones de Compute Optimal, el Token Efficiency mejora aproximadamente 2 veces: equivale a que aprender con 1 porción de datos sea comparable a que Adam aprenda con 2 porciones
    \item \textbf{Mayor dispersión}: al incorporar más parámetros mediante una mayor dispersión, cuantos más parámetros haya, mejor será la absorción de los mismos datos
    \item \textbf{Rephrase de datos}: reescribir los datos de alta calidad para que la misma porción de datos se aprenda varias veces en formas distintas, lo cual mejora la generalización
\end{enumerate}

\begin{knowledgebox}{Antecedentes técnicos del optimizador Muon}
El optimizador Muon fue propuesto por Keller. Su idea central es no considerar cada elemento de parámetro de la matriz de forma independiente (como hace Adam), sino tomar en cuenta la dependencia (dependency) entre los parámetros, y lograr un aprendizaje más eficiente mediante la ortogonalización de matrices.

La contribución de Moonshot AI fue llevar Muon a la práctica en modelos de lenguaje a gran escala. En una etapa temprana, el proyecto Moonlight lo validó por primera vez en un modelo de lenguaje de cierta escala; más adelante, al hacer Scale, surgieron nuevos problemas (como la explosión de MaxLogit), que se resolvieron proponiendo una nueva forma de Clipping.
\end{knowledgebox}

\begin{warningbox}{Reescribir datos no es lo mismo que copiarlos simplemente}
Aprender simplemente repitiendo la misma porción de datos provoca sobreajuste y problemas de generalización. La clave de la reescritura es introducir «nueva entropía»: hacer que los datos reescritos conserven el conocimiento central y, al mismo tiempo, tengan cierta diversidad, lo cual mejora la capacidad de generalización. Pero la elección de la forma de reescritura es crucial; cómo introducir nueva entropía de información sigue siendo una línea de investigación activa.
\end{warningbox}

\subsubsection{Dirección dos: capacidad Agentic y generalización}

\begin{importantbox}{La generalización de los Agent es el mayor desafío actual}
Para los modelos Agentic, el mayor desafío actual es la \textbf{generalización}:
\begin{itemize}
    \item Limitación de las técnicas de RL actuales: las tareas de entrenamiento y las métricas de evaluación suelen ser puntuales (por ejemplo, entrenar solo con SWE-Bench)
    \item Una mejora en la métrica no significa una mejora en la generalización: que la puntuación de SWE-Bench suba no implica que el modelo también sea mejor en otras tareas de código
    \item El problema de generalización de los Agent es más grave que en los modelos de conversación
    \item Los Benchmark son insuficientes o ya no son válidos
\end{itemize}
\end{importantbox}

Yang Zhilin (杨植麟) considera que una posible solución es «usar IA para entrenar IA»: hacer que el modelo participe en su propio proceso de entrenamiento y alineación, para mejorar la generalización de una manera más AI-Native.

\subsection{El proceso de investigación y desarrollo de K2}

\begin{itemize}
    \item La acumulación técnica comenzó un año antes (tecnologías como el optimizador Muon requieren experimentos de Scaling de ciclo largo para su validación)
    \item El ciclo de entrenamiento del propio K2 no fue largo, pero la preparación técnica previa requirió mucho tiempo
    \item El principal desafío encontrado: el optimizador Muon presentó el problema de explosión de MaxLogit durante el entrenamiento a gran escala (que no podía reproducirse en experimentos a pequeña escala)
    \item El equipo de investigación y el equipo de entrenamiento son el mismo equipo, porque los problemas durante el entrenamiento requieren una comprensión profunda de la tecnología subyacente para resolverse
\end{itemize}

\subsection{Resumen de la sección}

K1.5 validó la ruta de aprendizaje por refuerzo con Reward de extremo a extremo. K2 se centra en dos direcciones: (1) Token Efficiency (optimizador Muon + dispersión + reescritura de datos), (2) capacidad de generalización Agentic. La generalización de los Agent es el mayor cuello de botella actual, y quizá se necesite usar IA para entrenar IA para superarlo.

